\chapter{系统架构与工程实现}
\label{chap:architecture-implementation}

\section{总体架构设计}
\label{sec:impl-overall-architecture}

本文系统采用前后端分离和服务化架构。前端负责用户交互、阶段流程、图像展示、检测进度和报告页面；后端按照职责拆分为 Core API、Core BIZ、Activity Classification、Activity Reasoning、Activity Summary 和 ICW Common 等模块。基础设施包括 MySQL、Redis、MinIO 和 RocketMQ。这样的拆分并不是为了追求形式上的微服务，而是为了将不同性质的工作隔离：HTTP 接入、业务状态、Agent 调用、模型推理、报告生成和共享协议分别由不同模块承担。

Core API 是前端入口，提供 HTTP 和 WebSocket 能力。Core BIZ 是业务中心，负责数据库事务、对象存储、任务 Worker、回调处理和事件发布。Activity Classification 封装分类 Agent，Activity Reasoning 封装 Python 检测模型，Activity Summary 封装图像总结和项目报告 Agent。ICW Common 统一提供 Protocol Buffers 接口、枚举定义、gRPC 调用封装、日志工具、错误结构和通用工具函数。

这种架构的核心原则是“智能能力外置，业务状态内聚”。Agent 和模型能力被放在 Activity 服务中，可以独立演进；项目、图像、任务和报告状态集中在 BIZ 服务中，避免状态分散。API 服务不承担复杂业务，只负责把前端请求转换为 BIZ 调用，并把 BIZ 产生的事件推送给前端。这样的职责划分使系统既能接入外部智能能力，又能保持核心业务一致性。

从工作量角度看，本文系统不是简单调用一次大语言模型，而是需要把多个异构能力编排在同一业务流程中。Classification 调用的是外部 Agent 平台，Reasoning 调用的是本地 Python 模型，Summary 再次调用外部 Agent 平台，三类能力的运行时间、失败模式和返回格式均不同。系统必须通过统一协议和状态机把它们组织起来。

\begin{table}[!htb]
  \caption{系统服务划分}\label{tab:service-division}
  \centering
  \begin{tabular}{@{}p{0.22\linewidth}p{0.68\linewidth}@{}} \toprule
    \textbf{服务或模块} & \textbf{主要职责} \\ \midrule
    Core Web & 项目流程页面、图像资产管理、检测进度可视化、复核与报告展示。 \\
    Core API & HTTP 接口、鉴权入口、WebSocket 连接、RocketMQ 消息消费与前端广播。 \\
    Core BIZ & 业务状态机、数据库事务、对象存储封装、检测 Worker、回调落库和事件生产。 \\
    Activity Classification & 调用分类 Agent，输出需要执行的检测任务代码。 \\
    Activity Reasoning & 调用五类原子检测模型，上传检测产物并回调检测结果。 \\
    Activity Summary & 调用总结 Agent，生成单图总结和项目级评估报告。 \\
    ICW Common & 共享协议、枚举、gRPC 基础设施、日志、错误和工具函数。 \\ \bottomrule
  \end{tabular}
\end{table}

\section{通信与协议设计}
\label{sec:impl-communication}

系统外部通信采用 HTTP 和 WebSocket，内部服务通信采用 gRPC。前端通过 HTTP 发起项目、图像、检测、复核和报告相关请求；API 服务再通过 gRPC 调用 BIZ 服务。BIZ 调用 Activity 服务启动分类、推理和总结任务，Activity 服务执行完成后通过 gRPC 回调 BIZ。BIZ 将状态变化写入数据库并发布 RocketMQ 事件，API 服务消费事件后通过 WebSocket 推送给前端。

gRPC 的引入使服务间接口具有明确契约。系统中的任务状态、项目进度、检测任务代码、复核结论和报告状态均通过共享枚举定义，避免不同服务各自维护字符串常量。请求 ID 通过 gRPC metadata 透传，使 HTTP 请求、BIZ 处理、Activity 调用和回调日志能够关联起来。这对于复杂链路调试非常重要。

协议设计还承担了约束 Agent 输出的作用。分类结果、推理结果和总结结果都通过固定 gRPC 回调接口进入 BIZ。这样，即使 Agent 或模型服务内部实现发生变化，上层业务仍然只依赖稳定协议。接口定义统一放在 Common 模块中，并通过代码生成供 API、BIZ 和 Activity 服务使用，减少手写结构体转换带来的错误。

在接口设计中，本文尽量避免让调用方依赖被调用方内部结构。例如，Reasoning 服务不要求 BIZ 预先知道污渍和平整度任务会生成多少区域图像；BIZ 只提供受限上传前缀，Reasoning 执行后回调实际产物清单。再如，Summary 服务不直接读取数据库，而是接收 BIZ 生成的数据源或压缩上下文。这样，每个服务只暴露必要边界，内部实现可以独立变化。

\section{共享协议与枚举治理}
\label{sec:impl-common-protocol}

由于系统跨越 API、BIZ、Classification、Reasoning 和 Summary 多个 Go 仓库，如果每个仓库都自行定义请求结构、状态枚举和错误常量，很容易出现字段名不一致、枚举值不一致和转换函数膨胀等问题。本文将公共协议收敛到 ICW Common 模块，并通过 Protocol Buffers 生成代码。所有服务都依赖生成后的结构体和枚举，减少手写 DTO 转换。

共享协议治理在项目中体现为三个方面。第一，业务状态使用统一枚举，例如项目进度、图像状态、检测任务状态、子任务状态、报告状态等。第二，RPC 接口由 proto 文件定义，服务端实现接口，客户端调用生成代码。第三，API 层返回给前端的结构也尽量复用生成结构，减少 API 与 BIZ 之间重复定义。这样做提高了开发成本的前期投入，但降低了后续维护成本。

在 Agent 场景中，枚举治理尤其重要。分类 Agent 的输出最终会被解析为检测任务代码，如果代码在不同服务中以字符串散落维护，新增或修改任务时很容易遗漏。通过共享枚举和解析函数，系统可以在编译阶段暴露部分问题，并在运行时统一校验非法值。

\section{数据库与状态机实现}
\label{sec:impl-db-state}

数据库设计围绕项目、图像资产和检测任务展开。项目表保存建筑基础信息和项目进度；图像组和图像表保存资产组织关系；检测主任务表保存每张图像的检测状态；五张子任务表分别保存不同原子检测能力的结构化结果；图像总结表保存单图总结；项目报告表保存项目级报告结果。该设计虽然表数量较多，但可以将不同任务的专属指标清晰拆分，避免所有结果塞入一个无法维护的大字段。

检测主任务状态机是系统复杂度的集中体现。主任务从 pending 开始，经过 classifying、detecting、summarizing，最终进入 succeeded 或 failed。分类结果决定是否创建子任务，子任务结果决定是否进入总结，总结结果决定主任务是否完成。所有状态推进均由 BIZ 控制，Activity 服务只上报结果。这种设计避免外部模型服务直接操作核心数据库，提高系统一致性。

五类子任务表采用“通用字段 + 专属指标”的设计。通用字段包括任务 UUID、主任务 ID、用户 ID、项目 ID、图像 ID、状态、开始时间、完成时间和运行耗时；专属指标则根据检测任务差异定义。例如锈蚀检测需要记录锈蚀区域数量、锈蚀像素和锈蚀占比；平整度检测需要记录平整、不平整或非玻璃结论；爆裂检测只需要记录是否爆裂和置信度。相比所有结果统一存 JSON，这种设计便于后续查询、筛选和前端结构化展示。

数据库事务用于保证关键状态推进的原子性。例如，分类回调成功后，系统需要同时更新主任务路由字段、创建子任务、关联子任务 ID 并推进主任务到 detecting。如果这些操作分散执行，可能出现主任务已进入 detecting 但子任务尚未创建的中间状态。类似地，子任务回调时，系统需要在同一事务中写入检测结果、更新子任务状态、检查同一主任务下其他子任务状态，并在全部成功后创建总结任务。这些事务边界是系统一致性的核心。

数据库还承担了重试能力的基础。虽然本文当前阶段不展开单个子任务重试，但表结构和任务 UUID 已经为重试预留空间。失败任务可以通过主任务和子任务记录定位历史结果，后续可按图像维度删除旧子任务、清理对象存储并重新创建任务。没有结构化任务表，重试只能依赖文件状态或日志，难以成为可靠业务能力。

\section{Detection Worker 实现}
\label{sec:impl-worker}

Detection Worker 是检测流程的编排核心。用户启动检测时，BIZ 服务先为项目下图像创建主任务并返回，Worker 随后按配置的最大并发数消费任务。每个任务执行时，Worker 先将主任务推进到分类阶段，调用 Classification 服务；分类回调成功后，Worker 根据任务代码创建子任务并调用 Reasoning 服务；所有需要执行的子任务成功后，Worker 创建图像总结任务并调用 Summary 服务。

Worker 设计解决了同步调用难以支撑长任务的问题。如果前端请求一直等待所有模型执行结束，用户体验和服务稳定性都会受到影响。通过 Worker，启动检测接口只负责创建任务和入队，真正的检测链路在后台推进。前端通过 WebSocket 获得状态变化，从而实现“请求立即返回、任务持续执行、进度实时可见”的体验。

Worker 还负责处理状态之间的因果关系。例如，分类失败时不应继续创建子任务；分类成功但没有任何检测能力需要执行时，应直接完成主任务；任一必要子任务失败时，主任务最终应进入失败；全部子任务成功后，才能创建图像总结任务。这些逻辑如果分散在各 Activity 服务中，会导致状态难以维护。集中在 Worker 中处理，可以保持业务规则一致。

Worker 的另一个职责是并发控制。检测任务可能同时包含多张图像，如果不限制并发，分类 Agent、Python 模型、对象存储和数据库都会承受突发压力。系统通过最大并发配置限制后台任务推进速度，使吞吐量和资源占用保持平衡。并发控制不是简单性能参数，而是保证外部 Agent 平台、模型服务和存储服务稳定运行的必要条件。

Worker 与回调之间形成非阻塞协作关系。Worker 调用 Activity 服务时，只要求对方成功接收任务；真正执行结果通过回调返回。这种方式避免 BIZ Worker 长时间阻塞等待模型执行，也使 Activity 服务可以在内部使用协程和并发队列。对于外部 Agent 调用和深度学习推理这种长耗时任务，回调式架构比同步阻塞更适合。

\section{对象存储与产物管理}
\label{sec:impl-object-storage}

幕墙图像和检测产物体积较大，不适合直接存入数据库。系统使用 MinIO 作为对象存储，数据库只保存元数据和对象 Key。图像上传时，BIZ 生成预签名上传 URL，前端直接上传对象存储。检测产物上传时，由于污渍和平整度任务可能产生不定数量的区域图像，系统采用受限前缀 POST Policy，允许 Reasoning 服务在指定前缀下上传动态数量的 PNG 产物。

对象访问 URL 具有有效期，系统使用 Redis 缓存预签名下载 URL，降低重复生成成本。删除对象时，系统先删除 Redis 缓存，再删除 MinIO 对象，避免前端使用已失效缓存。对于检测产物，Reasoning 服务还计算 SHA256 并回调给 BIZ 保存，用于后续校验和追踪。

产物管理的难点在于不同检测任务的产物数量不同。裂缝检测通常输出原图、掩码图和浮层图；污渍和平整度检测可能根据区域数量输出多个编号图像；爆裂检测可能只有原图和报告。系统不能要求 BIZ 在调用前预知所有产物名称，因此采用前缀授权和回调产物清单的方式。Reasoning 服务实际发现并上传哪些文件，就在回调中告诉 BIZ。这样既支持动态产物数量，也避免业务服务和模型脚本强耦合。

对象存储设计还影响前端展示。检测结果页面需要根据任务类型展示不同产物组合：锈蚀任务展示标注图，裂缝任务展示浮层图和掩码图，污渍任务展示总体标注图以及每个区域的裁剪图和浮层图，平整度任务展示区域、线条、梯度和频域图。由于产物最终都以对象方式存储，前端只需要通过 BIZ 返回的预签名访问地址展示图像，而不需要理解对象存储内部权限机制。

\section{Reasoning 服务实现}
\label{sec:impl-reasoning}

Reasoning 服务是 Go 与 Python 模型能力的连接层。Go 服务负责 gRPC 接口、请求校验、运行目录准备、原图下载、Python Worker 调用、产物上传和回调。Python Worker 负责加载和执行具体检测模块。为避免每次检测都重新启动 Python 脚本和加载模型，系统采用常驻 Worker 和懒加载策略：某类模型首次调用时加载，后续任务复用已加载模型。

该实现兼顾了 Go 的工程能力和 Python 的模型生态。Go 适合处理网络、并发、对象存储和服务通信；Python 适合调用深度学习模型和图像处理库。通过明确边界，系统避免在 Python 中实现业务上传和数据库逻辑，也避免在 Go 中直接处理复杂模型推理。

Reasoning 服务的实现经历了从脚本执行到常驻 Worker 的优化。若每次任务都通过命令行启动 Python 脚本，则模型加载成本会重复发生，检测延迟较高。采用常驻 Python Worker 后，Go 服务通过标准输入输出或内部协议向 Worker 发送任务，Worker 在首次使用某个模型时加载并缓存。这样既保留 Python 模型实现，又降低重复加载成本。

运行目录设计也体现了工程约束。每次任务按任务代码和图像 UUID 创建独立目录，执行前清空旧目录并重新下载原图，避免重试或重复执行时读取历史产物。Python 模型只需要在约定目录中生成约定文件，Go 服务负责读取、上传和清理。这种目录约定比共享临时文件名更可靠，也便于调试和定位。

\section{前端可视化实现}
\label{sec:impl-frontend}

前端按照项目阶段组织页面。图像资产阶段提供上传、分组、批量操作和原图查看；智能检测阶段提供开始检测、失败重试、状态统计、实时进度和检测结果查看；人工复核阶段复用检测结果视图，并增加评论和准确性选择；报告阶段展示项目级评估报告。

检测进度可视化由两部分组成。一部分是图像卡片上的整体进度条，用于快速展示每张图像当前状态；另一部分是进度弹窗中的流程图，用于展示分类、并行原子检测、总结和结束节点。检测完成后，结果弹窗展示原始信息、子任务结果和图像总结。对于不同检测任务，系统根据产物数量自适应展示单图、双图或四图布局。

前端状态合并也是实现难点之一。WebSocket 消息可能因为网络、刷新或并发而出现延迟，前端不能简单用最后一条消息覆盖全部状态。系统因此同时提供状态查询接口和增量消息合并逻辑。页面首次进入或刷新时通过查询接口获得全量状态，运行过程中再用 WebSocket 消息更新对应图像和节点。这种设计提高了实时展示的可靠性。

检测结果展示采用任务类型驱动布局。不同原子任务的数据字段和图像产物不同，前端需要将数据库中的结构化结果转化为用户可读界面。例如，裂缝和污渍需要展示区域列表，平整度需要展示平整、不平整或非玻璃结果，爆裂任务没有区域数据。前端并非只显示 JSON，而是将字段翻译为中文指标，并根据区域序号切换对应产物。这样的展示方式使复杂检测数据能够被普通用户理解。

人工复核阶段复用检测结果视图，并在弹窗中增加评论和准确性选择。该设计避免重复开发两套结果页面，也保持用户体验一致。报告阶段则从检测图像视图转向项目报告视图，体现流程从图像级事实逐步走向项目级结论。

\section{部署实现}
\label{sec:impl-deployment}

系统最终以 Docker 镜像和 Docker Compose 方式部署。Go 服务镜像尽量保持轻量，Reasoning 服务由于需要 Python 环境和模型文件，采用 Go+Python 运行镜像并通过挂载目录提供模型权重。部署顺序上，先启动 RocketMQ，再初始化 Topic，随后启动三个 Activity 服务，再启动 BIZ，最后启动 API。这样的顺序保证依赖服务可用后，上层服务再开始接收请求。

部署过程中还需要处理服务地址、消息队列 NameServer、对象存储端点、Redis、MySQL 和 Agent 平台 Token 等环境配置。由于系统服务较多，部署配置本身也是工程工作量的一部分。本文通过显式环境变量和统一 Docker Compose 文件降低部署复杂度。

部署实践中还需要处理 RocketMQ Topic 初始化、Reasoning 模型文件挂载、Docker 镜像体积和服务启动顺序等问题。Reasoning 服务包含模型依赖，不适合作为纯 Go scratch 镜像；若把模型权重直接打入多架构镜像，镜像体积和推送稳定性都会受到影响。因此本文采用运行镜像与模型目录分离的方式，在服务器上通过只读挂载提供模型文件。该方案更适合大模型文件和检测权重的长期维护。

\section{工程复杂度总结}
\label{sec:impl-complexity-summary}

从表面看，系统使用的是成熟工程技术；但从整体链路看，工作量集中在多服务协同和状态一致性上。用户一次点击“开始检测”，背后会触发主任务创建、Worker 入队、分类 Agent 调用、分类回调、子任务创建、多个推理任务并发执行、产物上传、子任务结果落库、图像总结、报告准备、WebSocket 推送和前端状态合并等操作。每个环节都需要处理失败、重试、日志和数据一致性。

因此，本文工程实现不是简单页面和接口堆叠，而是为了支撑 Agent 自主调度和分层信息整合而构建的复杂运行环境。该环境使 Agent 能力从一次性调用转化为可持续执行的业务流程。

\begin{table}[!htb]
  \caption{一次图像检测任务涉及的主要工程环节}\label{tab:detection-engineering-steps}
  \centering
  \begin{tabular}{@{}p{0.26\linewidth}p{0.64\linewidth}@{}} \toprule
    \textbf{环节} & \textbf{说明} \\ \midrule
    任务创建 & 为图像创建主任务，记录用户、项目、图像和初始状态。 \\
    分类路由 & 调用分类 Agent，解析任务代码并校验输出合法性。 \\
    子任务生成 & 根据路由结果创建对应子任务，并关联到主任务。 \\
    原子检测 & 调用 Reasoning 服务执行模型，生成报告和图像产物。 \\
    产物上传 & 使用对象存储授权上传动态数量产物，并记录产物清单。 \\
    结果落库 & 将不同任务的结构化字段写入对应子任务表。 \\
    状态推进 & 判断子任务聚合状态，决定是否进入图像总结或失败终态。 \\
    实时通知 & 通过 RocketMQ 和 WebSocket 将状态变化推送给前端。 \\
    图像总结 & 调用 Summary 服务生成单图总结并写入数据库。 \\ \bottomrule
  \end{tabular}
\end{table}

\cref{tab:detection-engineering-steps} 展示了单张图像检测背后的主要工程环节。对于一个包含多张图像的项目，上述链路会并发执行多次，因此系统必须在并发控制、事务处理和前端状态合并上保持一致。这也是本文工程工作量的重要体现。
